% ------------------------------------------------------------
% References (verified volumes/pages/DOI -- see Notion draft v0.3)
% ------------------------------------------------------------
\begin{thebibliography}{14}

\bibitem{chen2017}
H.~Chen, Y.~Zhang, M.~K. Kalra, F.~Lin, Y.~Chen, P.~Liao, J.~Zhou, and G.~Wang,
``Low-dose CT with a residual encoder-decoder convolutional neural network,''
\emph{IEEE Transactions on Medical Imaging}, vol.~36, no.~12, pp. 2524--2535, 2017.
doi:10.1109/TMI.2017.2715284.

\bibitem{eulig2024}
E.~Eulig, B.~Ommer, and M.~Kachelrie{\ss},
``Benchmarking deep learning-based low-dose CT image denoising algorithms,''
\emph{Medical Physics}, vol.~51, no.~12, pp. 8776--8788, 2024.
doi:10.1002/mp.17379 (arXiv:2401.04661).

\bibitem{mccollough2017}
C.~H. McCollough et~al.,
``Low-dose CT for the detection and classification of metastatic liver lesions: Results of the 2016 Low Dose CT Grand Challenge,''
\emph{Medical Physics}, vol.~44, no.~10, pp. e339--e352, 2017.
doi:10.1002/mp.12345.

\bibitem{yang2018}
Q.~Yang, P.~Yan, Y.~Zhang, H.~Yu, Y.~Shi, X.~Mou, M.~K. Kalra, Y.~Zhang, L.~Sun, and G.~Wang,
``Low-dose CT image denoising using a generative adversarial network with Wasserstein distance and perceptual loss,''
\emph{IEEE Transactions on Medical Imaging}, vol.~37, no.~6, pp. 1348--1357, 2018.
doi:10.1109/TMI.2018.2827462.

\bibitem{huang2022}
Z.~Huang, J.~Zhang, Y.~Zhang, and H.~Shan,
``DU-GAN: Generative adversarial networks with dual-domain U-Net-based discriminators for low-dose CT denoising,''
\emph{IEEE Transactions on Instrumentation and Measurement}, vol.~71, pp. 1--12, 2022.
doi:10.1109/TIM.2021.3128703.

\bibitem{li2025}
L.~Li, W.~Wei, L.~Yang, W.~Zhang, J.~Dong, Y.~Liu, H.~Huang, and W.~Zhao,
``CT-Mamba: A hybrid convolutional state space model for low-dose CT denoising,''
\emph{Computerized Medical Imaging and Graphics}, vol.~124, art. 102595, 2025
(arXiv:2411.07930).

\bibitem{perez2018}
E.~Perez, F.~Strub, H.~de~Vries, V.~Dumoulin, and A.~Courville,
``FiLM: Visual reasoning with a general conditioning layer,''
in \emph{Proc. AAAI Conference on Artificial Intelligence}, vol.~32, no.~1, 2018.
doi:10.1609/aaai.v32i1.11671.

\bibitem{cui2025}
Y.~Cui, S.~W. Zamir, S.~Khan, A.~Knoll, M.~Shah, and F.~S. Khan,
``AdaIR: Adaptive all-in-one image restoration via frequency mining and modulation,''
in \emph{Proc. International Conference on Learning Representations (ICLR)}, 2025
(arXiv:2403.14614).

\bibitem{solomon2012}
J.~B. Solomon, O.~Christianson, and E.~Samei,
``Quantitative comparison of noise texture across CT scanners from different manufacturers,''
\emph{Medical Physics}, vol.~39, no.~10, pp. 6048--6055, 2012.
doi:10.1118/1.4752209.

\bibitem{he2016}
K.~He, X.~Zhang, S.~Ren, and J.~Sun,
``Deep residual learning for image recognition,''
in \emph{Proc. IEEE Conference on Computer Vision and Pattern Recognition (CVPR)}, pp. 770--778, 2016
--- ResNet trunk as implemented in the ldct-benchmark suite~\cite{eulig2024}.

\bibitem{solomon2020}
J.~Solomon, P.~Lyu, D.~Marin, and E.~Samei,
``Noise and spatial resolution properties of a commercially available deep learning-based CT reconstruction algorithm,''
\emph{Medical Physics}, vol.~47, 2020.
doi:10.1002/mp.14319.

\bibitem{zhang2025}
Z.~Zhang et~al.,
``Preserving diagnostic details in low-dose CT with frequency-domain guided deep learning,''
in \emph{Proc. Medical Image Computing and Computer-Assisted Intervention (MICCAI)}, LNCS, Springer, 2025.
doi:10.1007/978-3-032-09513-8\_48.

\bibitem{zhao2026}
Z.~Zhao, Y.~Wang, S.~Tian, X.~Li, and W.~Zhao,
``NIFA: Low-dose CT imaging via noise intensity field aware networks,''
\emph{Medical Image Analysis}, vol.~108, art. 103866, 2026.
doi:10.1016/j.media.2025.103866.

\bibitem{xu2025}
H.~Xu, L.~Chen, Y.~Zhang, G.~Quan, Y.~Xi, Y.~Chen, and X.~Ji,
``Task-based quantitative evaluation of single- or dual-domain networks for low dose CT,''
\emph{Biomedical Signal Processing and Control}, vol.~99, art. 106909, 2025.
doi:10.1016/j.bspc.2024.106909.

\end{thebibliography}
